\begin{figure}[H]
  \centering
  \begin{tikzpicture}[
    font=\small\sffamily,
    arr/.style={-{Latex}, thick},
    biarr/.style={{Latex}-{Latex}, thick},
    caja/.style={
      draw, thick, rounded corners=3pt, align=center,
      minimum width=3.5cm, minimum height=1.7cm,
      inner sep=5pt, fill=gray!6
    },
    nucleo/.style={
      draw, thick, rounded corners=3pt, align=center,
      minimum width=3.5cm, minimum height=1.7cm,
      inner sep=5pt, fill=gray!14
    },
    etiq/.style={font=\scriptsize\sffamily}
  ]
    % --- Profesional ---
    \node[inner sep=0pt, minimum width=0.8cm, minimum height=1.2cm]
      (act) {};
    \begin{scope}[shift={(act.center)}]
      \fill (0,0.40) circle (0.12);
      \draw[thick, line cap=round]
        (0,0.28) -- (0,-0.12)
        (-0.22,0.14) -- (0.22,0.14)
        (0,-0.12) -- (-0.17,-0.46)
        (0,-0.12) -- (0.17,-0.46);
    \end{scope}
    \node[below=0.10cm of act, align=center, font=\small\bfseries\sffamily]
      (actlbl) {Profesional};

    % --- Navegador y motor ---
    \node[caja, right=2.35cm of act] (nav)
      {\textbf{Navegador}\\[3pt]
       aplicación web};
    \node[nucleo, right=2.15cm of nav] (mot)
      {\textbf{Motor de reglas}\\[3pt]
       STOPP/START};

    % --- PDF (documento con esquina doblada) ---
    \node[
      draw, thick, align=center,
      minimum width=2.7cm, minimum height=1.85cm,
      inner sep=4pt, fill=gray!6,
      below=1.55cm of nav
    ] (pdf) {\textbf{PDF}\\[3pt]informe del caso};
    \fill[white] (pdf.north east) -- ++(-0.36,0) -- ++(0.36,-0.36) -- cycle;
    \draw[thick] ([xshift=-0.36cm]pdf.north east) |- ([yshift=-0.36cm]pdf.north east);

    \draw[arr] (act) -- node[above, etiq] {captura el caso} (nav);
    \draw[biarr] (nav) -- node[above, etiq] {evalúa / avisos} (mot);
    \draw[arr] (nav) -- node[right, etiq, pos=0.45] {exporta} (pdf);
  \end{tikzpicture}
  \caption{Visión de alto nivel: el profesional trabaja en el navegador;
  el motor de reglas evalúa el caso; el informe se exporta en PDF.}
  \label{fig:arq-cajas}
\end{figure}
